\section{A polynomial theorem for sparse active columns}
\label{sec:sparse}

\subsection{Replacing peeling by directed reachability}
Call a column \emph{inert} if none of its entries is one.  Every row
restriction still has no unit entry in that column, so it can never supply
a unit-column pivot.  Its sums on stopping sets can never equal one.
It may therefore be ignored for this decision, while its one-handle is
retained in the topological count.  Assume every non-inert column of $A$
has at most two nonzero rows.  Counts can still be arbitrarily large;
they are not assumed to be zero/one.  Define a
directed graph $P$ on the row set $R$ as follows.
\begin{itemize}
\item A column supported only on $r$ with $a_{rc}=1$ makes $r$ an initial
  seed.  A singleton entry of two or more makes no seed.
\item For a column supported on $\{u,v\}$, add an arc $v\to u$ if
  $a_{uc}=1$, and an arc $u\to v$ if $a_{vc}=1$.
\item Inert columns, including zero columns, make no contribution.
  Multiple arcs can be merged.
\end{itemize}
The direction records an implication: deleting the tail makes the head a
unit-column row.  Let $Z$ be the set of initial seeds.

\begin{lemma}[Exact reachability model]\label{lem:reachability}
For any set $Q$ of rows deleted by top cancellations, the union of $Q$
and all rows subsequently removed by unit-column peeling is exactly the
set reachable in $P$ from $Z\cup Q$.
\end{lemma}
\begin{proof}
Every initial seed can peel.  If an arc $v\to u$ comes from a column,
then that column has exactly one occurrence on $u$ and no possible other
row except $v$.  Once $v$ is deleted, $u$ can peel.  Induction on path
length therefore shows that every reachable row is removed.  The
certifying column cannot have been used and removed while its head row
was still live: at a unit-column pivot it had no other live owner.

Conversely, a row peeled at any step is certified either by a singleton
unit column, making it an initial seed, or by a two-row column whose other
row has already been removed.  In the latter case the corresponding
directed arc gives a reachability step.  Induction on the peel sequence
proves the reverse inclusion.
\end{proof}

Let the strongly connected components of $P$ be contracted into its
acyclic condensation graph.  Call a component a \emph{mandatory component}
if it is a source of this condensation graph and contains no initial seed.
List these components as $S_1,\ldots,S_q$.

\begin{lemma}[The only mandatory groups]\label{lem:mandatory}
Deleting $Q$ allows complete peeling if and only if
\[
                        Q\cap S_i\ne\varnothing
                           \quad(1\leq i\leq q).
\]
The $S_i$ are pairwise disjoint stopping sets.
\end{lemma}
\begin{proof}
A source component with no seed cannot be reached from another component.
It therefore needs a row of $Q$.  Conversely every component of a finite
acyclic graph is reachable from some source.  If each source either contains
a seed or is met by $Q$, all rows are reachable, and
Lemma~\ref{lem:reachability} applies.

To see that a mandatory component is stopping, consider a column having
an occurrence on a row of $S_i$.  If it has only one row and count one,
that row would be a seed.  If it has exactly one row in $S_i$, another row
outside $S_i$, and count one on the inside row, it would give an incoming
arc to the source component.  Both cases are impossible.  If both rows
are inside, the total is at least two; if the only inside count is larger
than one, the same is true.  An inert column can never sum to one.
Hence every column has zero or at least two
occurrences on $S_i$.
\end{proof}

\subsection{Independent representatives in a graph}
Each mandatory component is a set of row labels and therefore a set of
edges of $D$, with ineligible rows represented by root loops.  The groups
are disjoint.  We need one edge from each group,
with all chosen edges independent in the graphic matroid.  Such a partial
forest can be extended to a spanning tree, and any spanning tree meeting
the groups contains such a partial forest.

For completeness we state and prove the independent-representatives
criterion, a classical result of Rado \cite{Rado1942}.

\begin{lemma}[Rado's rank criterion]\label{lem:rado}
Let $M$ be a matroid and $S_1,\ldots,S_q$ finite subsets of its ground set.
There exist distinct representatives $e_i\in S_i$ forming an independent
set if and only if
\begin{equation}\label{eq:rado}
       \rk_M\left(\bigcup_{i\in J}S_i\right)\geq |J|
       \qquad\text{for every }J\subseteq\{1,\ldots,q\}.
\end{equation}
\end{lemma}
\begin{proof}
Necessity follows by restricting an independent representative set to $J$.
The case $q=0$ is empty.  For sufficiency induct on $q$.  The case $q=1$
says that $S_1$ contains
a nonloop, which follows from its positive rank.

Suppose a nonempty proper subfamily $J$ is tight in \eqref{eq:rado}.
By induction choose independent representatives $B$ for $J$.  They have
cardinality and rank $|J|$, so they span $U_J=\bigcup_{i\in J}S_i$.
For a set $K$ of remaining groups, contraction by $B$ gives
\[
 \rk_{M/B}(U_K)
   =\rk_M(U_K\cup B)-|J|
   =\rk_M(U_K\cup U_J)-|J|
   \geq |K|.
\]
Apply induction in the contracted matroid to the remaining groups.  Their
representatives together with $B$ give the required independent set.

If there is no nonempty proper tight subfamily, every such subfamily has
rank at least its size plus one.  Choose any nonloop $e\in S_q$.
For a nonempty set $K\subseteq\{1,\ldots,q-1\}$,
\[
            \rk_{M/e}(U_K)\geq\rk_M(U_K)-1\geq |K|.
\]
Induction after contracting $e$ supplies the remaining representatives;
adjoining $e$ finishes the proof.  The contraction treats $e$ as absent
from the remaining ground set, so the representatives are distinct.
\end{proof}

\begin{theorem}[Polynomial forest choice in the two-row regime]
\label{thm:degree-two}
Suppose $D$ is a connected rooted multigraph, and every non-inert column
of $A$ has at most two nonzero rows.  Compute the mandatory components
$S_1,\ldots,S_q$ of $P$.  Then the following are equivalent:
\begin{enumerate}
\item $(D,A)$ has a forest-compatible raw completion;
\item the groups $S_i$ have an independent system of representative edges
  in the graphic matroid of $D$;
\item for every set $J$ of mandatory groups,
\begin{equation}\label{eq:graphic-rado}
             \rk_D\left(\bigcup_{i\in J}S_i\right)\geq |J|.
\end{equation}
\end{enumerate}
These conditions can be decided in polynomial time.  A positive certificate
is a spanning tree with an explicit complete peel order.  A negative
certificate is a subfamily $J$ violating \eqref{eq:graphic-rado}, together
with the mandatory-component data, which can be recomputed independently.
\end{theorem}
\begin{proof}
Lemma~\ref{lem:mandatory} reduces completion to choosing a spanning tree
meeting each mandatory group.  Such a tree contains independent
representatives.  Conversely an independent representative set extends
to a spanning tree, which meets every group and therefore succeeds.
Lemma~\ref{lem:rado} proves the equivalence with the rank inequalities.

For the algorithmic assertion, restrict the ground set to the union of
the groups.  Intersect the graphic matroid with the partition matroid
allowing at most one edge from each group.  A common independent set
of size $q$ gives the representatives.  Matroid intersection has a
polynomial augmenting-path algorithm with independence oracles
\cite{Edmonds1979}.  Here the oracles are elementary: a
disjoint-set forest detects graphic cycles, and color counts test the
partition condition.  The residual exchange graph of a failed maximum
intersection produces a rank defect; equivalently the matroid-intersection
min--max theorem and Lemma~\ref{lem:rado} supply the subfamily $J$.
The delivered implementation extracts such a witness and its checker
verifies the inequality directly, without trusting an optimality assertion.
\end{proof}

\subsection{Explicit algorithm and complexity scope}
The algorithm first scans columns to detect the unsupported case of an
active column with more than two incident rows.  It then computes seeds,
arcs, strongly connected
components, and mandatory groups.  These steps take linear time in the
matrix incidence size, with standard adjacency-list graph algorithms.

Starting from the empty common independent set, the matroid-intersection
algorithm builds an exchange graph.  Its vertices are candidate edges.
Arcs encode exchanges that preserve one of the two matroids, and endpoints
encode single-element additions that preserve the other.  A shortest
augmenting path toggles membership and increases the common independent
set by one.  If no path exists, the reachable set yields the min--max
certificate.  There are at most $q\leq|R|$ augmentations.  Even a conservative
implementation using $O(|R|^2)$ oracle calls per augmentation and rebuilding
a graphic oracle in $O(|R|+|V(D)|)$ operations is polynomial:
\[
 O\bigl(I+|R|^3(|R|+|V(D)|)\bigr),
\]
where $I$ is the occurrence-input size.  No enumeration of all subsets $J$
in \eqref{eq:graphic-rado} is required.  The implementation and measured
operation counts are documented in Sections~\ref{sec:implementation}
and~\ref{sec:experiments}.

The column restriction is a geometric crossing restriction in an
authenticated core, not a bound on algebraic rank or on the number of
vertices of $D$.  Large counts on either of the two rows are permitted.
For example, an abstract solver input column $(1,2^{1000})^{\mathsf T}$
contributes just the arc
from the second row to the first.  A column with entries $(2,2,2)$ is
inert and harmless for the restriction.  A column with entries $(1,1,1)$
is outside the theorem.  These large-count examples concern the generalized
finite solver, not the original native columns bounded below.

\subsection{Why shortest augmentation works, and where the obstruction comes from}
Here is a self-contained justification for the specialized exchange step.
Restrict to the union $G$ of the mandatory groups and maintain a rainbow
forest $F$.  The exchange graph has sources $e\notin F$ for which $F+e$
is a forest and targets $e\notin F$ of an unrepresented group.  For
$f\in F$ and $e\notin F$ it has arcs
\[
 f\longrightarrow e\quad\hbox{if }F-f+e\hbox{ is a forest},\qquad
 e\longrightarrow f\quad\hbox{if }F-f+e\hbox{ is rainbow}.
\]
Take a shortest source-to-target path
$e_0,f_1,e_1,\ldots,f_\ell,e_\ell$.  None of its intermediate outside
vertices is a target, so $e_{i-1}$ has the group of $f_i$, whereas
$e_\ell$ has a new group.  The symmetric difference is therefore rainbow
and has $|F|+1$ edges.

For graphic independence, represent edges by vertex-incidence vectors
over $\mathbb F_2$.  The forest vectors are independent.  The source
$e_0$ is outside their span.  Every later $e_i$ is not a source, since
otherwise its suffix would shorten the path.  It therefore has a unique
fundamental cycle in $F$.  Its expansion in the forest basis has
coefficient one at $f_i$ and coefficient zero at each $f_j$ with $j<i$;
a nonzero earlier coefficient would supply the shortcut $f_j\to e_i$.
Modulo $F\setminus\{f_1,\ldots,f_\ell\}$ these expansions form a
triangular matrix with unit diagonal.  They are independent, and adding
$e_0$, outside the original span, preserves independence.  The augmented
set is a larger rainbow forest.

If no target is reachable, let $W$ be the reachable set of edge labels.
Every group appearing in $W$ is represented by an edge of $F\cap W$:
an unrepresented reachable group would give a target, and an outside
edge in a represented group has an arc to that group's selected edge.
Thus exactly $|F\cap W|$ groups appear in $W$.  Each outside edge not
in $W$ is spanned by $F\setminus W$: it is not a source, and its
fundamental cycle has no reachable selected edge, which would give an
incoming exchange arc.  Hence
\[
                   \rk_D(G\setminus W)=|F\setminus W|.
\]
Let $J$ consist of the groups with no edge in $W$.  If $|F|<q$, then
\[
 \rk_D\left(\bigcup_{j\in J}S_j\right)
 \leq |F\setminus W|
 =|F|-|F\cap W|<q-|F\cap W|=|J|.
\]
This supplies the claimed negative certificate directly.  The checker
only needs the stopping sets and the rank inequality, not the internal
exchange graph.

\subsection{A geometric bound isolating the remaining native case}
\begin{proposition}[Three occurrences per native cocore]\label{prop:three}
For the unsigned matrix of original native core two-handles against
native one-handle cocores,
\[
                         \sum_r a_{rc}\leq3
                    \qquad\hbox{for every column }c.
\]
The same holds after contraction of a fixed lower zero-/one-handle forest
and any dual-tree row deletion.  It is not asserted after adding boundary
relations or performing later handle slides.
\end{proposition}
\begin{proof}
Work in one abstract triangular source face before quotient identifications.
Its disjoint normal arcs have endpoints on distinct local sides.  Every
complementary disk region meets each of the three sides in at most one
gap.  This follows inductively: an additional arc splits one region,
divides its intervals on the two endpoint sides between the children,
and leaves each other side interval wholly in at most one child.
Initially the assertion holds for the uncut triangle.

The cocore of a native one-handle is one such face region.  Each geometric
traversal by a core two-handle corresponds to an actual boundary-gap
occurrence on that cocore, of which there are at most three.  Different
occurrences remain counted even if quotient identifications give them
the same global edge or handle label.  Removing parallelity two-handles
can only decrease the count.  Lower forest contraction deletes its
columns and introduces no crossings of retained cocores.  Dual-tree
normalization deletes rows while retaining their attaching charts.
\end{proof}

Consequently, up to permutation, every nonzero native column has
positive-entry multiset
\[
        (1),\ (2),\ (3),\ (1,1),\ (1,2),\ \hbox{or }(1,1,1).
\]
The last shape is the only one outside Theorem~\ref{thm:degree-two}.
For its three rows $u,v,w$, peeling gives the two-premise implications
$\{u,v\}\to w$, $\{u,w\}\to v$, and $\{v,w\}\to u$.
This is a precise point at which ordinary directed reachability ceases
to suffice.

\begin{proposition}[Parity compatibility with the dual graph]
\label{prop:cycle}
Use the actual unsaturated source counts in this statement.  Before top
cancellations, the dual edges whose entries in a fixed source column are
odd form an even-degree subgraph of $D$.  The same statement
holds for the retained rows in $D/Q$ after contracting a selected dual
forest $Q$.
\end{proposition}
\begin{proof}
Work over $\mathbb F_2$.  In the handle chain complex, the boundary of
a three-handle records the actual cap incidences of its two-handles
modulo two.  Thus $\partial_3$ is the incidence matrix of $D$ with the
formal root row omitted, transposed to the handle-boundary convention.
The corresponding row of $\partial_2$ is the unsigned source column
modulo two.  The identity $\partial_2\partial_3=0$ implies even degree
at each three-handle vertex.  Loops contribute twice.  The root also
has even degree, since the sum of all graph degrees is even.  Contraction
of a nonloop combines its endpoint parity equations and deletes that
edge; an even subgraph remains even in the contracted multigraph.
Iterate along $Q$.
\end{proof}

A native $(1,1,1)$ column is therefore a three-edge mod-two cycle:
a triangle, a parallel pair together with a loop, or three loops.
A singleton odd entry must label a dual loop.  These are necessary
geometric compatibility conditions, not a realization theorem for arbitrary
matrices.  The nonmatroid example in Section~\ref{sec:forests} already
has column totals $(2,3,3)$.  It can also satisfy the parity condition:
make $r_0,r_2,r_3$ three parallel root-to-vertex edges and $r_1$ a loop.
Its three columns are then even subgraphs.  Thus neither of these two
necessary conditions restores the matroid exchange axiom.  No native
triangulated realization of that example is claimed.

\subsection{The boundary between polynomial and general search}
Theorem~\ref{thm:degree-two} does not establish an intrinsic complexity
dichotomy for all handle structures.  We do not prove the unrestricted
problem NP-hard, and an improved algorithm may exist for other classes.
The exact obstruction-guided procedure remains available when the column
restriction fails.  Its canonical residual and graphic contraction can
also expose a smaller two-row instance later in the search.

A practical implementation can therefore use the polynomial solver first,
then use the general search with an explicit budget.  Positive output
always carries a tree and checked peel order.  Exhaustion under a budget
must produce an inconclusive status; it must not be confused with a
negative rank certificate or an exhaustively checked search tree.
